% Generated draft: requires mathematical and source review.
\documentclass[11pt]{article}
\usepackage[margin=25mm]{geometry}
\usepackage{fontspec}
\setmainfont{Latin Modern Roman}
\usepackage{amsmath,amssymb,mathtools,unicode-math}
\setmathfont{Latin Modern Math}
\usepackage{enumitem,needspace,xurl,hyperref}
\setlength{\parindent}{0pt}
\setlength{\parskip}{5pt}
\setlength{\emergencystretch}{3em}
\newcommand{\N}{\mathbb{N}}
\newcommand{\Z}{\mathbb{Z}}
\newcommand{\Q}{\mathbb{Q}}
\newcommand{\R}{\mathbb{R}}
\newcommand{\C}{\mathbb{C}}
\newcommand{\F}{\mathbb{F}}
\newcommand{\A}{\mathbb{A}}
\newcommand{\PP}{\mathbb{P}}
\newcommand{\E}{\mathbb{E}}
\newcommand{\eps}{\varepsilon}
\newcommand{\dd}{\,\mathrm{d}}
\begin{document}
\section*{UnsolvedMath ID: 20003274}
\textbf{Source problem code:} catalogue entry
\subsection{Definitions and mathematical statement}
L. Silberman's source request is to extend unipotent-flow measure rigidity to the adelic setting. It does not specify the global field, adelic acting subgroup, quotient, or whether full infinite-adic factors are involved. The aid separates characteristic-zero finite-\(S\) classification, attributed to Margulis--Tomanov and Ratner, from a general full-adelic or parabolic formulation, which is not automatically covered.

Its precise compact-tail statement is the following. Let commuting groups \(L,K\) act continuously on a suitable locally compact second-countable space \(X\), with \(K\) compact and a well-behaved orbit quotient \(q:X\to X/K\). Finite \(L\)- and \(K\)-invariant measures on \(X\) correspond bijectively by \(q_*\) to finite \(L\)-invariant measures \(\nu\) on \(X/K\). The inverse uses normalized Haar orbital measures:
\[
 \int_X\varphi\,\dd\mu
 =\int_{X/K}\left(\int_K\varphi(xk)\,\dd m_K(k)\right)\dd\nu([x]).
\]
The inner expression is independent of the orbit representative; no free-action hypothesis is needed when stabilizers are handled by the Haar orbit measure. The aid states
\[
 \mu\text{ is jointly }LK\text{-ergodic}
 \Longleftrightarrow\nu\text{ is }L\text{-ergodic}.
\]
It does not replace joint ergodicity by \(L\)-ergodicity upstairs.

For an adelic application over \(\Q\), let \(Y=\mathbf G(\Q)\backslash\mathbf G(\mathbb A)\), let \(S\) be finite including the archimedean place, and choose a maximal-compact tail \(K^S\) outside \(S\). Assume the requisite finite-double-class and arithmetic hypotheses so that
\[
 Y/K^S\cong\coprod_{j=1}^{r}\Gamma_j\backslash G_S
\]
is a finite disjoint union of \(S\)-arithmetic lattice quotients. If \(L\subseteq G_S\) is in the characteristic-zero unipotent-generated classification setting, every finite measure invariant under both \(L\) and the \emph{full} \(K^S\) is recovered without loss from its finite-level measure, and the finite-\(S\) classification applies to the latter. The exact algebraic-group, lattice, and unipotent-generation hypotheses of the cited classification are not supplied.

The aid emphasizes a proper-tail obstruction. Already on a product \(Z\times K\), with \(L\) acting on \(Z\) and a proper compact subgroup \(U\subsetneq K\) acting on the second factor, the product of an \(L\)-ergodic measure on \(Z\) and Haar measure on one \(U\)-coset is jointly \(LU\)-ergodic and has the same projection to \(Z\), but is not \(K\)-invariant. Thus invariance under a compact unipotent tail does not, on those assumptions alone, imply maximal-compact tail invariance.

It also works with the upper unipotent subgroup \(\mathbf U\subset\operatorname{SL}_2\), whose closed orbit in \(\operatorname{SL}_2(\Q)\backslash\operatorname{SL}_2(\mathbb A)\) is
\[
 \mathbf U(\Q)\backslash\mathbf U(\mathbb A)\cong\Q\backslash\mathbb A.
\]
This compact additive quotient has normalized Haar probability and illustrates full adelic unipotent invariance without silently replacing it by maximal-compact tail invariance. The detailed orbit calculation and chosen tail are not fully supplied in the record.

The original trailing discussion expects algebraically manageable finite ergodic measures for other parabolic-like systems. For connected semisimple Lie \(H\subseteq G\) and parabolic \(P\subseteq H\), it reports that finite \(P\)-invariant measures on \(G/\Gamma\) are \(H\)-invariant [146]. The precise quotient and theorem hypotheses require checking; this historical Lie-group assertion is not itself a proved infinite-adelic analogue. The aid's compact-tail reduction leaves the missing general full-adelic step separate.
\subsection{Short English statement}
The supplied compact-tail theorem gives an exact finite-level reduction only for measures invariant under the whole compact tail. Smaller unipotent-tail invariance does not imply that hypothesis. Finite-\(S\) classification and the original broader adelic/parabolic request therefore remain distinct.
\subsection{Sources}
\begin{itemize}
\item Local UnsolvedMath registry, permanent ID \(20003274\), exact source code \texttt{catalogue entry}; canonical \texttt{aim-other-notes.json}, \texttt{notes[80]}, AIM domain Other, registry category Miscellaneous, source tag \texttt{problem}.
\item AIM workshop ``Emerging applications of measure rigidity,'' Problem 7 attributed to L. Silberman: \url{https://aimath.org/WWN/measrigid/measrigid.pdf}. Statement status \texttt{reconstructed\_unverified}; the supplied verification distinguishes the original grammatical defect from line-break and spacing artifacts. No PDF was fetched.
\item Original source-local label [146] concerns the Lie parabolic-invariance assertion. The aid names Margulis--Tomanov and Ratner for characteristic-zero finite-\(S\) unipotent measure classification, but supplies no full authoritative citations or complete theorem hypotheses. The original literature field is empty.
\item Research attempt 1: result \texttt{partial\_result}, problem status at run \texttt{partially\_solved}, rigor label \texttt{proved}, novelty confidence \texttt{low}; classification \texttt{PARTIAL-PROGRESS}, registry status \texttt{partially\_solved}. Created and updated 14 August 2026. These are supplied assessments requiring review.
\item Attribution: Ulam AI / UnsolvedMath Contributors, CC BY 4.0. The machine-generated research aid is not peer reviewed and requires independent verification before citation.
\end{itemize}
\end{document}
